\begin{proof}[Proof of \cref{obj:lem:hidden-allocation-contraction}]
We use the convention \(0^0=1\) and interpret empty products as one.

\begin{enumerate}
\item First establish the scalar series estimate that will close the argument. Define
\[
e:=\exp(1),
\qquad
\rho_{\nu_{\mathrm{rev}},\nu_{\mathrm{hid}}}
:=
\min\left\{1,\frac{4(\nu_{\mathrm{rev}}+\nu_{\mathrm{hid}})}{B}\right\}
\quad
(\nu_{\mathrm{rev}},\nu_{\mathrm{hid}}\in\mathbb N_0),
\]
and
\[
\mathfrak w_{\nu_{\mathrm{rev}},\nu_{\mathrm{hid}}}
:=
\binom B{\nu_{\mathrm{rev}}}
\binom{B-\nu_{\mathrm{rev}}}{\nu_{\mathrm{hid}}}
\rho_{\nu_{\mathrm{rev}},\nu_{\mathrm{hid}}}^{\nu_{\mathrm{hid}}}
(p\eta_1)^{\nu_{\mathrm{rev}}}\eta^{\nu_{\mathrm{hid}}}
\quad
(0\le\nu_{\mathrm{rev}},\nu_{\mathrm{hid}}\le B).
\]
By \cref{obj:lem:walsh-channel-energy},
\(0\le\eta\le\eta_1\). Also \(0\le(1-q)^d\le1\), so \(0\le p\le1\).
Consequently all these terms are nonnegative.

Define the nonconstant sum by
\[
\mathfrak S_{\mathrm{nc}}
:=
\sum_{\substack{0\le\nu_{\mathrm{rev}},\nu_{\mathrm{hid}}\le B\\
\nu_{\mathrm{rev}}+\nu_{\mathrm{hid}}>0}}
\mathfrak w_{\nu_{\mathrm{rev}},\nu_{\mathrm{hid}}}
=
\sum_{\nu_{\mathrm{rev}}=0}^B
\sum_{\nu_{\mathrm{hid}}=0}^B
\mathfrak w_{\nu_{\mathrm{rev}},\nu_{\mathrm{hid}}}-1.
\]
The equality follows because \(\mathfrak w_{0,0}=1\).

The falling-factorial formula for binomial coefficients gives
\[
\binom B{\nu_{\mathrm{rev}}}
\le\frac{B^{\nu_{\mathrm{rev}}}}{\nu_{\mathrm{rev}}!},
\qquad
\binom{B-\nu_{\mathrm{rev}}}{\nu_{\mathrm{hid}}}
\le\frac{B^{\nu_{\mathrm{hid}}}}{\nu_{\mathrm{hid}}!}.
\]
Using \(B>0\) and the definition of \(\rho\), we obtain
\[
\begin{aligned}
&\binom B{\nu_{\mathrm{rev}}}
\binom{B-\nu_{\mathrm{rev}}}{\nu_{\mathrm{hid}}}
\rho_{\nu_{\mathrm{rev}},\nu_{\mathrm{hid}}}^{\nu_{\mathrm{hid}}}
\\
&\qquad\le
\frac{B^{\nu_{\mathrm{rev}}}}{\nu_{\mathrm{rev}}!}
4^{\nu_{\mathrm{hid}}}
\frac{(\nu_{\mathrm{rev}}+\nu_{\mathrm{hid}})^{\nu_{\mathrm{hid}}}}
{\nu_{\mathrm{hid}}!}
\\
&\qquad\le
\frac{B^{\nu_{\mathrm{rev}}}}{\nu_{\mathrm{rev}}!}
4^{\nu_{\mathrm{hid}}}
\exp(\nu_{\mathrm{rev}}+\nu_{\mathrm{hid}})
=
\frac{(eB)^{\nu_{\mathrm{rev}}}}{\nu_{\mathrm{rev}}!}
(4e)^{\nu_{\mathrm{hid}}}.
\end{aligned}
\]
The second inequality bounds one nonnegative term of the exponential series by its sum. It includes \(\nu_{\mathrm{hid}}=0\), with the stated power convention. Thus
\[
\begin{aligned}
\mathfrak S_{\mathrm{nc}}+1
&\le
\left(
\sum_{\nu_{\mathrm{rev}}=0}^B
\frac{(eBp\eta_1)^{\nu_{\mathrm{rev}}}}{\nu_{\mathrm{rev}}!}
\right)
\left(
\sum_{\nu_{\mathrm{hid}}=0}^B
(4e\eta)^{\nu_{\mathrm{hid}}}
\right)
\\
&\le
\exp(eBp\eta_1)\sum_{\nu_{\mathrm{hid}}=0}^{\infty}
(4e\eta)^{\nu_{\mathrm{hid}}}
=
\frac{\exp(eBp\eta_1)}{1-4e\eta}.
\end{aligned}
\]
Here \(0\le4e\eta<1\) justifies the geometric series. Therefore
\[
\mathfrak S_{\mathrm{nc}}
\le
\frac{\exp(eBp\eta_1)}{1-4e\eta}-1.
\]
% lean: CausalSmith.Experimentation.ThinnedgraphAdditiveRiskFrontier.contraction_nonconstant_sum_bound

\item We next identify the design and its conditional reference. Define
\[
\mu_Z:=\operatorname{Law}(Z),
\qquad
\mu_{\mathrm{ass}}
:=
\bigotimes_{j\in V}\operatorname{Bernoulli}(1/2),
\qquad
\mu_{\mathrm{aud}}
:=
\bigotimes_{\substack{(j,i)\in V^2\\j\ne i}}
\operatorname{Bernoulli}(q).
\]
The audit measure is a probability measure because \(q\in[0,1]\).
The first marginal of \(\mu_Z\otimes\mu_{\mathrm{aud}}\) is therefore
\(\mu_Z\). The assignment condition in
\cref{obj:ass:assignment}, together with assignment and audit independence, gives
\[
\mu_Z=\mu_{\mathrm{ass}},
\qquad
\operatorname{Law}(Z,W)
=
\mu_{\mathrm{ass}}\otimes\mu_{\mathrm{aud}}.
\]
The source partition is drawn independently of this design, and the retained graph is a function of the partition and audit marks. Hence
\[
\operatorname{Law}(H,Z)
=
\operatorname{Law}(H)\otimes\mu_{\mathrm{ass}}.
\]
In particular, the assignment signs remain independent uniform signs conditional on every graph atom of positive mass.

Define the response reference measure on \(\mathbb R^B\) by
\[
\nu_{\mathrm{resp}}(dy)
:=
\prod_{\ell=1}^B g_d(y_\ell)\,dy_\ell.
\]
Each \(g_d\) is a probability density, being a finite mixture of translates of the baseline density in \cref{obj:lem:block-law}. Thus
\[
Q_H=\mu_{\mathrm{ass}}\otimes\nu_{\mathrm{resp}}.
\]
Every realized retained graph is compatible with its sampled source partition. Positive graph atoms therefore admit compatible partitions, and they exhaust the graph marginal almost surely. All graphwise constructions below concern these atoms; auxiliary graph functions are extended by zero at null atoms.
% lean: CausalSmith.Experimentation.ThinnedgraphAdditiveRiskFrontier.hidden_allocation_contraction

\item Fix such a graph \(H\). Let the fixed source population, the revealed source sets, and the hidden source set be
\[
\mathcal S_{\mathrm{src}}:=\bigcup_{\ell=1}^B S_\ell,
\qquad
R_\ell:=\{j\in\mathcal S_{\mathrm{src}}:
\exists i\in C_\ell,\ (j,i)\in H\},
\qquad
\mathcal D:=\mathcal S_{\mathrm{src}}\setminus\bigcup_{\ell=1}^B R_\ell.
\]
The union defining \(\mathcal S_{\mathrm{src}}\) is the same fixed set for every source partition. Compatibility makes the sets \(R_\ell\) pairwise disjoint. Write
\[
r_\ell:=|R_\ell|,
\qquad
u_\ell=d-r_\ell,
\qquad
M=\sum_{\ell=1}^B u_\ell=|\mathcal D|.
\]
For \(z\in\{0,1\}^V\), define
\[
X_j(z):=2z_j-1,
\qquad
X_j:=X_j(Z),
\qquad
A_\ell(z):=\sum_{j\in R_\ell}X_j(z),
\qquad
K(z):=\sum_{j\in\mathcal D}z_j.
\]

The family of hidden-label completions is
\[
\mathfrak C_H
:=
\left\{
(\mathcal D_1,\ldots,\mathcal D_B):
\mathcal D=\bigsqcup_{\ell=1}^B\mathcal D_\ell,\
|\mathcal D_\ell|=u_\ell
\right\}.
\]
It is nonempty and has cardinality
\[
|\mathfrak C_H|=\frac{M!}{\prod_{\ell=1}^B u_\ell!}.
\]
For integers \(0\le k_\ell\le u_\ell\) with
\(\sum_\ell k_\ell=K(z)\), counting the allocations of treated and untreated hidden labels separately gives
\[
\begin{aligned}
&\frac{
\left|
\left\{
(\mathcal D_\ell)\in\mathfrak C_H:
|\mathcal D_\ell\cap\{j:z_j=1\}|=k_\ell
\text{ for every }\ell
\right\}
\right|
}{|\mathfrak C_H|}
\\
&\qquad=
\frac{
\displaystyle
\frac{K(z)!}{\prod_\ell k_\ell!}
\frac{(M-K(z))!}{\prod_\ell(u_\ell-k_\ell)!}
}{
\displaystyle
\frac{M!}{\prod_\ell u_\ell!}
}
=
\frac{\prod_\ell\binom{u_\ell}{k_\ell}}{\binom M{K(z)}}.
\end{aligned}
\]
The count is zero when the total-degree condition fails. Since \(0\le K(z)\le M\), its denominator is positive, including \(M=K(z)=0\).

Expanding the coefficient-extraction formula in
\cref{obj:lem:block-law} and applying this count gives the exact identity
\[
\lambda_{+1,h}(y\mid H,z)
=
\frac1{|\mathfrak C_H|}
\sum_{(\mathcal D_\ell)\in\mathfrak C_H}
\prod_{\ell=1}^B
f\left(
y_\ell-\frac h{2d}
\left(A_\ell(z)+\sum_{j\in\mathcal D_\ell}X_j(z)\right)
\right).
\]
Thus the conditional density is the uniform-completion average of the original labeled row densities.

If \(g_d(y_\ell)=0\), every translated density contributing to its nonnegative finite mixture is zero. The last display then gives
\(\lambda_{+1,h}(y\mid H,z)=0\). Consequently a version of the likelihood is
\[
L_+(z,y)
=
\begin{cases}
\displaystyle
\frac{\lambda_{+1,h}(y\mid H,z)}
{\prod_{\ell=1}^B g_d(y_\ell)},
&\prod_{\ell=1}^B g_d(y_\ell)>0,\\[6pt]
0,&\prod_{\ell=1}^B g_d(y_\ell)=0.
\end{cases}
\]
The response reference assigns full measure to the first case. The conditional density identification and the product form of \(Q_H\) therefore yield
\[
\chi_H^2
=
\int_{\mathbb R^B}
\int_{\{0,1\}^V}
\bigl(L_+(z,y)-1\bigr)^2
\,\mu_{\mathrm{ass}}(dz)\,\nu_{\mathrm{resp}}(dy).
\]
Finite channel expansions are bounded, so the integral rearrangements below are justified. Multiplying this identity by the remaining-capacity indicator preserves it almost surely under the actual graph marginal.
% lean: CausalSmith.Experimentation.ThinnedgraphAdditiveRiskFrontier.hiddenChiSq_good_integral_eq_actualDensityEnergy

\item We now express this density energy through labeled coefficients. For
\(0\le k\le M\), define
\[
\Psi_k(X_{\mathcal D})
:=
\binom Mk^{-1}
\sum_{\substack{J_{\mathrm{hid}}\subseteq\mathcal D\\
|J_{\mathrm{hid}}|=k}}
\prod_{j\in J_{\mathrm{hid}}}X_j,
\qquad
X_{\mathcal D}:=(X_j)_{j\in\mathcal D}.
\]
Every denominator is positive; in particular \(\Psi_0=1\).

For a hidden-degree vector, define
\[
\boldsymbol k:=(k_1,\ldots,k_B)
\in\prod_{\ell=1}^B\{0,\ldots,u_\ell\},
\qquad
|\boldsymbol k|:=\sum_{\ell=1}^B k_\ell,
\qquad
\omega(\boldsymbol k):=\prod_{\ell=1}^B\binom{u_\ell}{k_\ell}.
\]
For any fixed completion \((\mathcal D_\ell)\), averaging over all permutations of the \(M\) hidden labels gives
\[
\begin{aligned}
&\frac1{M!}\sum_{\pi\in\operatorname{Sym}(\mathcal D)}
\sum_{\substack{
J_{\mathrm{hid},\ell}\subseteq\mathcal D_\ell\\
|J_{\mathrm{hid},\ell}|=k_\ell\ \text{for every }\ell}}
\prod_{\ell=1}^B
\prod_{j\in J_{\mathrm{hid},\ell}}X_{\pi(j)}
\\
&\qquad=
\omega(\boldsymbol k)\,
\Psi_{|\boldsymbol k|}(X_{\mathcal D}).
\end{aligned}
\]
Indeed, there are \(\omega(\boldsymbol k)\) selections. For each selection, its union has size \(|\boldsymbol k|\), and a uniform permutation maps that union uniformly onto subsets of this size. The identity also holds for \(M=0\), when the permutation and selection families are singletons. Uniform completion averaging is invariant under these permutations.

For each revealed subset \(J\subseteq\bigcup_\ell R_\ell\), define
\[
r_{J,\ell}:=|J\cap R_\ell|,
\]
and, for \(0\le k\le M\) and \(y\in\mathbb R^B\), define
\[
\mathcal C_{J,k}(y)
:=
\sum_{\substack{
\boldsymbol k\in\prod_\ell\{0,\ldots,u_\ell\}\\
|\boldsymbol k|=k}}
\omega(\boldsymbol k)
\prod_{\ell=1}^B a_{r_{J,\ell}+k_\ell}(y_\ell).
\]
The orders in this formula lie between zero and \(d\).
Applying the row expansion in \cref{obj:lem:walsh-channel-energy} to the completion average, and then grouping by total hidden degree, gives
\[
L_+(Z,y)
=
\sum_{J\subseteq\bigcup_\ell R_\ell}
\left(\prod_{j\in J}X_j\right)
\sum_{k=0}^M
\Psi_k(X_{\mathcal D})\,\mathcal C_{J,k}(y)
\]
for \(\nu_{\mathrm{resp}}\)-almost every \(y\). The disjointness of the revealed row sets identifies a tuple of rowwise revealed subsets with its union \(J\).

Independent uniform signs satisfy
\[
\mathbb E\!\left[
\Psi_k(X_{\mathcal D})\Psi_{k'}(X_{\mathcal D})
\mid H
\right]
=
\frac{\mathbf1\{k=k'\}}{\binom Mk}
\qquad(0\le k,k'\le M).
\]
To see this, expand the two subset sums. The expected product is zero unless the two selected subsets coincide, since a sign appearing exactly once has mean zero. There are \(\binom Mk\) coincident pairs when \(k=k'\). The same argument for revealed subsets, and independence of the revealed and hidden signs, gives
\[
\begin{aligned}
&\mathbb E\!\left[
\left(\prod_{j\in J}X_j\right)\Psi_k(X_{\mathcal D})
\left(\prod_{j\in J'}X_j\right)\Psi_{k'}(X_{\mathcal D})
\mid H
\right]
\\
&\qquad=
\frac{\mathbf1\{J=J',\ k=k'\}}{\binom Mk}.
\end{aligned}
\]
Treatment coordinates outside the source population contribute a factor one.

For a square-integrable function
\(\varpi:\mathbb R^B\to\mathbb R\), write
\[
\|\varpi\|_{\mathrm{resp}}^2
:=
\int_{\mathbb R^B}\varpi(y)^2\,\nu_{\mathrm{resp}}(dy).
\]
The zero-degree coefficient is
\[
\mathcal C_{\varnothing,0}(y)=\prod_{\ell=1}^B a_0(y_\ell)=1
\quad\text{for }\nu_{\mathrm{resp}}\text{-almost every }y.
\]
Subtracting this coefficient, applying sign orthogonality, and integrating over \(y\) proves the exact labeled-energy identity
\[
\chi_H^2
=
\sum_{\substack{
J\subseteq\bigcup_\ell R_\ell,\ 0\le k\le M\\
(J,k)\ne(\varnothing,0)}}
\frac{\|\mathcal C_{J,k}\|_{\mathrm{resp}}^2}{\binom Mk}.
\]
% lean: CausalSmith.Experimentation.ThinnedgraphAdditiveRiskFrontier.actualDensityEnergy_eq_actualLabelEnergy

\item Separate these coefficients by their response supports. Define
\[
\mathcal B_J:=\{\ell:r_{J,\ell}>0\}.
\]
For \(\mathcal E\subseteq\{1,\ldots,B\}\setminus\mathcal B_J\) and
\(0\le k\le M\), let
\[
\mathscr K_{J,\mathcal E,k}
:=
\left\{
\boldsymbol k\in\prod_\ell\{0,\ldots,u_\ell\}:
|\boldsymbol k|=k,\
k_\ell>0\ (\ell\in\mathcal E),\
k_\ell=0\ (\ell\notin\mathcal B_J\cup\mathcal E)
\right\},
\]
and define
\[
\mathcal C_{J,k,\mathcal E}(y)
:=
\sum_{\boldsymbol k\in\mathscr K_{J,\mathcal E,k}}
\omega(\boldsymbol k)
\prod_{\ell=1}^B a_{r_{J,\ell}+k_\ell}(y_\ell).
\]
Every hidden-degree vector belongs to exactly one such family, so
\[
\mathcal C_{J,k}
=
\sum_{\mathcal E\subseteq\{1,\ldots,B\}\setminus\mathcal B_J}
\mathcal C_{J,k,\mathcal E}.
\]

Each summand has positive coefficient order precisely on
\(\mathcal B_J\cup\mathcal E\). If \(\mathcal E\ne\mathcal E'\), one row belongs to exactly one of these supports. At that row the cross product has one positive-order coefficient and one factor \(a_0=1\) almost everywhere. The zero-mean assertion of
\cref{obj:lem:walsh-channel-energy}, together with the product response measure, gives
\[
\int
\mathcal C_{J,k,\mathcal E}(y)
\mathcal C_{J,k,\mathcal E'}(y)\,
\nu_{\mathrm{resp}}(dy)=0.
\]
It follows that
\[
\|\mathcal C_{J,k}\|_{\mathrm{resp}}^2
=
\sum_{\mathcal E\subseteq\{1,\ldots,B\}\setminus\mathcal B_J}
\|\mathcal C_{J,k,\mathcal E}\|_{\mathrm{resp}}^2.
\]
% lean: CausalSmith.Experimentation.ThinnedgraphAdditiveRiskFrontier.actualLabelWalshCoeff_support_energy

\item Suppose now that \(M\ge Bd/4\). Since \(B,d\ge1\), this implies \(M>0\).
For fixed \(J,\mathcal E\), define
\[
\nu_{\mathrm{rev}}:=|\mathcal B_J|,
\qquad
\nu_{\mathrm{hid}}:=|\mathcal E|,
\qquad
U_{\mathrm{act}}:=\sum_{\ell\in\mathcal B_J\cup\mathcal E}u_\ell.
\]
These quantities satisfy
\[
U_{\mathrm{act}}\le M,
\qquad
U_{\mathrm{act}}\le d(\nu_{\mathrm{rev}}+\nu_{\mathrm{hid}}),
\qquad
\frac{U_{\mathrm{act}}}{M}
\le\rho_{\nu_{\mathrm{rev}},\nu_{\mathrm{hid}}}\le1.
\]
The first two inequalities use the nonnegative capacities and \(u_\ell\le d\); the third uses both \(U_{\mathrm{act}}\le M\) and \(M\ge Bd/4\).

Every vector in \(\mathscr K_{J,\mathcal E,k}\) has
\[
k=\sum_{\ell=1}^B k_\ell
\ge\sum_{\ell\in\mathcal E}k_\ell
\ge\nu_{\mathrm{hid}}.
\]
For \(0\le k\le U_{\mathrm{act}}\), the falling-factorial identity gives
\[
\frac{\binom{U_{\mathrm{act}}}{k}}{\binom Mk}
=
\prod_{v=0}^{k-1}
\frac{U_{\mathrm{act}}-v}{M-v}
\le
\left(\frac{U_{\mathrm{act}}}{M}\right)^k.
\]
Every denominator in this product is positive. At \(k=0\), both sides are one. If \(U_{\mathrm{act}}<k\le M\), the numerator binomial coefficient is zero. Consequently, whenever \(0\le k\le M\) and \(k\ge\nu_{\mathrm{hid}}\),
\[
\frac{\binom{U_{\mathrm{act}}}{k}}{\binom Mk}
\le
\rho_{\nu_{\mathrm{rev}},\nu_{\mathrm{hid}}}^{\nu_{\mathrm{hid}}}.
\]
Here the exponent decreases from \(k\) to \(\nu_{\mathrm{hid}}\) because the capacity fraction lies in \([0,1]\).
% lean: CausalSmith.Experimentation.ThinnedgraphAdditiveRiskFrontier.hidden_choose_ratio_good_event

The allocation weights satisfy
\[
\sum_{\boldsymbol k\in\mathscr K_{J,\mathcal E,k}}
\omega(\boldsymbol k)
\le
[x^k]\prod_{\ell\in\mathcal B_J\cup\mathcal E}(1+x)^{u_\ell}
=
\binom{U_{\mathrm{act}}}{k},
\]
because the restricted family is contained in the full family selecting \(k\) slots from these active rows.

Weighted Cauchy--Schwarz gives, pointwise in \(y\),
\[
\begin{aligned}
\mathcal C_{J,k,\mathcal E}(y)^2
&\le
\left(
\sum_{\boldsymbol k\in\mathscr K_{J,\mathcal E,k}}
\omega(\boldsymbol k)
\right)
\\
&\quad{}\times
\sum_{\boldsymbol k\in\mathscr K_{J,\mathcal E,k}}
\omega(\boldsymbol k)
\prod_{\ell=1}^B
a_{r_{J,\ell}+k_\ell}(y_\ell)^2.
\end{aligned}
\]
Extend the energy notation to degree zero by
\[
\gamma_0:=
\int_{\mathbb R}a_0(w)^2g_d(w)\,dw=1.
\]
This equality follows from \(a_0=1\) on positive reference support and \(\int g_d=1\).
Independence of the reference responses yields
\[
\int
\prod_{\ell=1}^B
a_{r_{J,\ell}+k_\ell}(y_\ell)^2\,
\nu_{\mathrm{resp}}(dy)
=
\prod_{\ell=1}^B\gamma_{r_{J,\ell}+k_\ell}.
\]
Thus
\[
\frac{\|\mathcal C_{J,k,\mathcal E}\|_{\mathrm{resp}}^2}{\binom Mk}
\le
\rho_{\nu_{\mathrm{rev}},\nu_{\mathrm{hid}}}^{\nu_{\mathrm{hid}}}
\sum_{\boldsymbol k\in\mathscr K_{J,\mathcal E,k}}
\omega(\boldsymbol k)
\prod_{\ell=1}^B\gamma_{r_{J,\ell}+k_\ell}.
\]
For \(k<\nu_{\mathrm{hid}}\), the restricted family is empty and both sides vanish.

Summing over \(k=0,\ldots,M\) covers each admissible degree vector exactly once. The remaining restrictions are rowwise, so the sum factors:
\[
\begin{aligned}
&\sum_{k=0}^M
\frac{\|\mathcal C_{J,k,\mathcal E}\|_{\mathrm{resp}}^2}{\binom Mk}
\\
&\quad\le
\rho_{\nu_{\mathrm{rev}},\nu_{\mathrm{hid}}}^{\nu_{\mathrm{hid}}}
\prod_{\ell\in\mathcal B_J}
\left(
\sum_{k_\ell=0}^{u_\ell}
\binom{u_\ell}{k_\ell}\gamma_{r_{J,\ell}+k_\ell}
\right)
\prod_{\ell\in\mathcal E}
\left(
\sum_{k_\ell=1}^{u_\ell}
\binom{u_\ell}{k_\ell}\gamma_{k_\ell}
\right).
\end{aligned}
\]
Define the active-support envelope by
\[
\begin{aligned}
\mathcal A_H
:=
\sum_{J\subseteq\bigcup_\ell R_\ell}
\sum_{\mathcal E\subseteq\{1,\ldots,B\}\setminus\mathcal B_J}
&\rho_{|\mathcal B_J|,|\mathcal E|}^{|\mathcal E|}
\prod_{\ell\in\mathcal B_J}
\left(
\sum_{k_\ell=0}^{u_\ell}
\binom{u_\ell}{k_\ell}\gamma_{r_{J,\ell}+k_\ell}
\right)
\\[-2pt]
&{}\times
\prod_{\ell\in\mathcal E}
\left(
\sum_{k_\ell=1}^{u_\ell}
\binom{u_\ell}{k_\ell}\gamma_{k_\ell}
\right).
\end{aligned}
\]
Restoring the coefficient \((J,k)=(\varnothing,0)\) adds exactly one to the labeled energy. The support decomposition and the last bound therefore give
\[
\chi_H^2+1\le\mathcal A_H,
\qquad
\chi_H^2\le\mathcal A_H-1
\quad\text{when }M\ge Bd/4.
\]
If \(M<Bd/4\), multiplying either quantity by the remaining-capacity indicator gives zero. This includes \(M=0\).
% lean: CausalSmith.Experimentation.ThinnedgraphAdditiveRiskFrontier.actualLabelEnergy_le_activeEnvelope

\item Reindex the envelope by revealed-active and hidden-active rows. For integers \(0\le r\le d\), define
\[
E_{\mathrm{rev}}(r)
:=
\sum_{s_{\mathrm{rev}}=1}^{r}
\binom r{s_{\mathrm{rev}}}
\sum_{k_{\mathrm{row}}=0}^{d-r}
\binom{d-r}{k_{\mathrm{row}}}
\gamma_{s_{\mathrm{rev}}+k_{\mathrm{row}}},
\]
and
\[
E_{\mathrm{hid}}(r)
:=
\sum_{k_{\mathrm{row}}=1}^{d-r}
\binom{d-r}{k_{\mathrm{row}}}\gamma_{k_{\mathrm{row}}}.
\]
The sums are empty when their upper limits are below their lower limits.

Since the \(R_\ell\) are disjoint, each global revealed subset \(J\) corresponds bijectively to its rowwise subsets \(J\cap R_\ell\).
For a row in \(\mathcal B_J\), summing over its nonempty revealed subsets gives
\[
\sum_{\substack{J_\ell\subseteq R_\ell\\J_\ell\ne\varnothing}}
\sum_{k_\ell=0}^{u_\ell}
\binom{u_\ell}{k_\ell}\gamma_{|J_\ell|+k_\ell}
=
E_{\mathrm{rev}}(r_\ell).
\]
For a hidden-active row, its factor is \(E_{\mathrm{hid}}(r_\ell)\).
It follows that
\[
\mathcal A_H
=
\sum_{\substack{
\mathcal B_{\mathrm{rev}},\mathcal B_{\mathrm{hid}}
\subseteq\{1,\ldots,B\}\\
\mathcal B_{\mathrm{rev}}\cap\mathcal B_{\mathrm{hid}}=\varnothing}}
\rho_{|\mathcal B_{\mathrm{rev}}|,|\mathcal B_{\mathrm{hid}}|}^{|\mathcal B_{\mathrm{hid}}|}
\prod_{\ell\in\mathcal B_{\mathrm{rev}}}E_{\mathrm{rev}}(r_\ell)
\prod_{\ell\in\mathcal B_{\mathrm{hid}}}E_{\mathrm{hid}}(r_\ell).
\]
The pair of empty row sets contributes exactly one. Removing it gives
\[
\begin{aligned}
\mathcal A_H-1
=
\sum_{\substack{
\mathcal B_{\mathrm{rev}}\cap\mathcal B_{\mathrm{hid}}=\varnothing\\
|\mathcal B_{\mathrm{rev}}|+|\mathcal B_{\mathrm{hid}}|>0}}
&\rho_{|\mathcal B_{\mathrm{rev}}|,|\mathcal B_{\mathrm{hid}}|}^{|\mathcal B_{\mathrm{hid}}|}
\\
&{}\times
\prod_{\ell\in\mathcal B_{\mathrm{rev}}}E_{\mathrm{rev}}(r_\ell)
\prod_{\ell\in\mathcal B_{\mathrm{hid}}}E_{\mathrm{hid}}(r_\ell).
\end{aligned}
\]
The row sets in this display are subsets of \(\{1,\ldots,B\}\).
% lean: CausalSmith.Experimentation.ThinnedgraphAdditiveRiskFrontier.actualLabelActiveEnvelope_sub_one_eq_row_energy

\item We bound the average of this nonconstant envelope under the detailed graph law. Counting a nonempty subset of the \(d\) row slots according to whether it contains a revealed slot gives
\[
\begin{aligned}
E_{\mathrm{rev}}(r)+E_{\mathrm{hid}}(r)
&=
\sum_{s=1}^d
\left(
\sum_{s_{\mathrm{rev}}=0}^s
\binom r{s_{\mathrm{rev}}}
\binom{d-r}{s-s_{\mathrm{rev}}}
\right)\gamma_s
\\
&=
\sum_{s=1}^d\binom ds\gamma_s=\eta.
\end{aligned}
\]
The second equality is the binomial convolution identity, with out-of-range binomial coefficients zero. All energies are nonnegative, so
\[
E_{\mathrm{rev}}(r)\ge0,
\qquad
0\le E_{\mathrm{hid}}(r)\le\eta.
\]

Fix a source partition \(S\). A source has \(d\) possible outgoing arrows, each audited independently. Its association-revelation indicator therefore has probability \(p=1-(1-q)^d\). Different sources use disjoint audit coordinates, so these indicators are independent conditional on \(S\), including across rows. For any specified \(s\)-subset of a row, the probability of at least one revealed member is \(1-(1-p)^s\). Thus
\[
\mathbb E\!\left[E_{\mathrm{rev}}(r_\ell)\mid S\right]
=
\sum_{s=1}^d
\binom ds\bigl(1-(1-p)^s\bigr)\gamma_s.
\]
The Bernoulli inequality gives
\[
(1-p)^s\ge1-sp,
\qquad
1-(1-p)^s\le sp
\quad(s\in\mathbb N_0).
\]
For completeness, the first inequality starts with equality at \(s=0\); multiplying it by \(1-p\ge0\) gives
\[
(1-p)^{s+1}
\ge(1-sp)(1-p)
=1-(s+1)p+sp^2
\ge1-(s+1)p.
\]
Consequently,
\[
\mathbb E\!\left[E_{\mathrm{rev}}(r_\ell)\mid S\right]
\le p\eta_1,
\qquad
\mathbb E\!\left[E_{\mathrm{hid}}(r_\ell)\mid S\right]
\le\eta.
\]

Every term in \(\mathcal A_H-1\) is nonnegative. On the remaining-capacity event its indicator leaves the term unchanged; outside that event the indicated term is zero. Hence
\[
0\le
\mathbf1\{M\ge Bd/4\}(\mathcal A_H-1)
\le\mathcal A_H-1.
\]
For disjoint row sets, conditional independence gives
\[
\begin{aligned}
&\mathbb E\!\left[
\prod_{\ell\in\mathcal B_{\mathrm{rev}}}E_{\mathrm{rev}}(r_\ell)
\prod_{\ell\in\mathcal B_{\mathrm{hid}}}E_{\mathrm{hid}}(r_\ell)
\;\middle|\;S
\right]
\\
&\quad=
\prod_{\ell\in\mathcal B_{\mathrm{rev}}}
\mathbb E[E_{\mathrm{rev}}(r_\ell)\mid S]
\prod_{\ell\in\mathcal B_{\mathrm{hid}}}
\mathbb E[E_{\mathrm{hid}}(r_\ell)\mid S]
\\
&\quad\le
(p\eta_1)^{|\mathcal B_{\mathrm{rev}}|}
\eta^{|\mathcal B_{\mathrm{hid}}|}.
\end{aligned}
\]
Empty row sets give empty products equal to one.

There are
\(\binom B{\nu_{\mathrm{rev}}}\) choices of
\(\mathcal B_{\mathrm{rev}}\) of size \(\nu_{\mathrm{rev}}\), followed by
\(\binom{B-\nu_{\mathrm{rev}}}{\nu_{\mathrm{hid}}}\) choices of a disjoint
\(\mathcal B_{\mathrm{hid}}\) of size \(\nu_{\mathrm{hid}}\).
When their sizes exceed \(B\) in total, the latter coefficient is zero.
Summing the preceding product bounds therefore gives, for every fixed partition,
\[
\mathbb E\!\left[
\mathbf1\{M\ge Bd/4\}(\mathcal A_H-1)
\mid S
\right]
\le\mathfrak S_{\mathrm{nc}}.
\]

The detailed retained graph generated by a partition and audit array is
\[
H(S,W)
:=
\{(j,i):
\exists\ell,\ j\in S_\ell,\ i\in C_\ell,\ W_{ji}=1\}.
\]
Averaging the fixed-partition inequality over the uniform partition law thus averages over the actual marginal of this full graph:
\[
\begin{aligned}
\mathbb E_H\!\left[
\mathbf1\{M\ge Bd/4\}(\mathcal A_H-1)
\right]
&=
\mathbb E_S\!\left[
\mathbb E\!\left[
\mathbf1\{M\ge Bd/4\}(\mathcal A_{H(S,W)}-1)
\mid S
\right]\right]
\\
&\le\mathfrak S_{\mathrm{nc}}.
\end{aligned}
\]
This average includes every source label and every retained-edge subset.
% lean: CausalSmith.Experimentation.ThinnedgraphAdditiveRiskFrontier.retained_row_energy_average_le

\item Combining the exact conditional density and labeled-energy identities with the good-event envelope bound yields
\[
\begin{aligned}
\mathbb E_H\!\left[
\mathbf1\{M\ge m/4\}\chi_H^2
\right]
&\le
\mathbb E_H\!\left[
\mathbf1\{M\ge Bd/4\}(\mathcal A_H-1)
\right]
\\
&\le\mathfrak S_{\mathrm{nc}}
\\
&\le
\frac{\exp(eBp\eta_1)}{1-4e\eta}-1,
\end{aligned}
\]
where \(m=Bd\). This is the asserted bound under the actual complete retained-graph marginal.
% lean: CausalSmith.Experimentation.ThinnedgraphAdditiveRiskFrontier.hidden_allocation_contraction
\end{enumerate}
\end{proof}
